
\begin{flushright}{\slshape
    Du musst bereit sein Dinge zu tun.} \\ \medskip
    --- A meme on the internet, 2022.
\end{flushright}



\bigskip

\begingroup
\let\clearpage\relax
\let\cleardoublepage\relax
\let\cleardoublepage\relax
\begin{otherlanguage}{ngerman}
\chapter*{Vorwort}
\pdfbookmark[1]{Vorwort}{vorwort}
Eine Promotion zu starten benötigt nicht viel: die vage Motivation etwas zum Wissensschatz der eigenen Disziplin beizutragen, sich über einen längeren Zeitraum intensiv in ein Thema einzugraben, und, vielleicht, auch einen Mangel an konkreten Alternativen.
Eine Promotion abzuschließen benötigt dagegen vor allem eines: Durchhaltevermögen. 

Als ich Anfang 2018 direkt nach dem Studium mit der Promotion anfing, hatte ich eigentlich keine große Vorstellung davon, was es eigentlich heißt, eine Dissertation zu verfassen. Ich hatte gerade eine aufwändige Masterarbeit geschrieben und verteidigt, mit allen Höhen und Tiefen die das so mit sich bringt, und  hätte eigentlich auch gut ein paar Wochen (oder Monate) Lust gehabt, erstmal nichts zu tun. Stattdessen hatte mein Betreuer (und zukünftiger Doktorvater) eine Stelle als wissenschaftlicher Mitarbeiter für die statistische Beratung zur Hand, welche ich so interessant fand, dass ich meine Erholungszeit auf eine Woche reduzierte. Interesse an einer Promotion hatte ich auch, allerdings ohne konkret zu wissen, zu welchem Thema (das Thema der Masterarbeit konnte ich nicht mehr sehen). 

Ich startete also in meine wissenschaftliche Karriere ohne großen Plan. In den ersten zwei Jahren beschäftigte ich mich lose mit verschiedenen Themen und griff zum Schluss ein Thema meines Doktorvaters auf. Das Thema --- nicht-asymptotische Konfidenzbereiche auf Mannigfaltigkeiten --- war ein eher theoretisches, eines in welchem Daten und statistische Ergebnisse eher Mittel zum mathematischen Zweck, denn Selbstzweck sind. Parallel betrieb ich die statistische Beratung, in der das Ziel genau umgekehrt schien: die statistischen Ergebnisse und daraus abgeleiteten Erkenntnisse waren das Ziel (manchmal auf die erfolgreiche Veröffentlichung oder Antworten auf die Fragen von Reviewern), die Methodik dagegen eher Mittel zum Zweck. Diese kognitive Dissonanz ging nicht lange gut --- für die Promotion musste ich tief in the Theorie abtauchen, für die Beratung hingegen angewandt breit aufgestellt sein, um auf die vielen unterschiedlichen Fragen gut eingehen zu können. 

Daher war es fast schon ein Glücksfall, dass die COVID-19 Pandemie im Frühjahr 2020 dafür sorgte, dass die Welt aus den Fugen geriet. Anfangs dachte ich noch, dass ich nun endlich Zeit haben würde mich um meine Forschung zu kümmern, aber mein Doktorvater hatte andere Ideen: als Statistiker könnten, ja sollten, wir doch etwas zur Erforschung der Pandemie beitragen. Anfangs etwas widerwillig (ich dachte ja, dass ich nun Zeit für meine Dissertation hätte), dann aber mit mehr und mehr Interesse begann ich also, mich mit mathematischer und angewandter Epidemiologie zu beschäftigen. Irgendwann (ich weiß nicht mehr genau wann) hatte ich die Erkenntnis, dass dieses Thema mir helfen würde, genau diese Dissonanz zwischen der angewandten Beratung und der theoretischen Forschung aufzuheben. Ich wechselte also nach zwei Jahren mein Promotionsthema --- und habe diesen Schritt bis heute nicht bereut. 

Durch die Pandemie lag natürlich viel Aufmerksamkeit auf meinem neu auserkorenen Thema, was dazu führte, dass wir schnell Teil einer größeren Gruppe von Statistiker:innen und Epidemiolog:innen wurden, die sich damit beschäftigten den Verlauf der Pandemie vorherzusagen. In wöchentlichen Treffen tauschten wir uns zu Vorhersagen, den neuesten epidemiologischen Erkenntnissen, Datenartefakten und mehr aus. Diesen wöchentlichen Treffen schreibe ich einen Großteil meines Durchhaltevermögens zu --- sie gaben mir das Gefühl, nicht alleine mit den Forschungsthemen da zu stehen, insbesondere in den manchmal schweren (und rückblickend absurden) Zeiten der Lockdowns. 

Zwischen dem Frühjahr 2020 und Mitte 2023 hatte ich das Vergnügen, verschiedenste Themen in der angewandten Epidemiologie zu bearbeiten und dabei auch einige Publikationen mit Kolleg:innen aus vielen Disziplinen zu veröffentlichen. Der Fokus dieser Veröffentlichungen, wie auch meine Beiträge zu diesen, war größtenteils angewandter Natur, in den kurzlebigen Zeiten der Epidemie (Lockdowns, Impfungen, verschiedene Varianten) lag mein Fokus nicht auf mathematischen Ergebnissen, stattdessen hangelte ich mich von Vorhersage zu Vorhersage (Fallzahlen für Deutschland und Polen, dann EU-weit, dann Hospitalisierungen für Deutschland). Am Ende des Tages strebte ich aber eine Promotion an einer mathematischen-naturwissenschaftlichen Fakultät an, ein bisschen Mathematik sollte also schon vorkommen. Glücklicherweise war ich an verschiedenen Stellen darauf gestoßen, dass Importance Sampling eine nützliche Technik im epidemiologischen Kontext sein könnte, und dass die Anwender:innen dieser Methoden sich gegenseitig nicht zu kennen schienen; jedenfalls zitierten sie sich nicht gegenseitig. 

Dadurch entstand langsam die Fragestellung, die heute den mathematischen Kern der Arbeit bildet: welche dieser beiden Methoden sollte man eigentlich benutzen? Empirisch schien die Cross-Entropy Methode deutlich instabiler als Efficient Importance Sampling, aber warum ist das so? Nach einem (leider erfolglosen) Abstecher, die Cross-Entropy Methode für inverse Probleme zu nutzen, machte ich mich also daran, diese Frage zu formalisieren und --- in Teilen --- zu beantworten. 

Leider ging mir so langsam die Zeit aus. Meine Finanzierung endete im Januar 2024 und ich hatte mir fest vorgenommen, meine Stelle nicht weiter zu verlängern, um meine Promotion endlich abzuschließen. Rückblickend bin ich froh, dass ich zu diesem Zeitpunkt nicht wusste, wie viel Arbeit noch vor mir lag. Ich hatte zwar einen guten Plan für die Inhalte meiner Dissertation, aber weder veröffentlichte Ergebnisse, die ich nutzen konnte, noch einen Großteil der Modelle für den Anwendungsteil --- diese wollte ich alle neu implementieren und aufziehen --- noch eine Vorstellung davon, ob die mathematischen Ergebnisse, die ich noch zu produzieren gedachte, interessant sein könnten. Anfang 2024 hatte ich also mit viel zu tun und wenig vorzuweisen, das Dissertationsdokument bestand damals gerade mal aus 40 Seiten. Ich fand relativ schnell einen Job in Leipzig und konnte dort einen späten Start zum Oktober vereinbaren, sodass ich ein halbes Jahr Zeit hatte, Ergebnisse zu schaffen und diese zu Papier zu bringen. 

Hätte ich ein bisschen mehr mathematische Vorarbeit gehabt, hätte das auch klappen können. So war es nicht möglich --- ich schrieb das große Mathematikkapitel mehrmals um, nachdem mir mehrfach Fehler in den Beweisen unterkamen. Auch die Struktur und Inhalte gingen durch mehrere Iterationen, genauso wie die angewandten Ergebnisse. Es stellte sich heraus, dass es tatsächlich schwierig war, neue Modelle aufzustellen, anzufitten und deren Ergebnisse anwendungsnah sinnvoll einzuordnen. Im September 2024 stand ich also mit einer halbfertigen Dissertationsschrift --- große Teile des Anwendungskapitels fehlten --- ohne Feedback --- mein Doktorvater hatte sich zwar immer wieder mit mir getroffen, aber noch keine Zeile Text gelesen --- und einem neuen Job, in dem ich auch durchstarten wollte, da. Die nächsten $1\frac12$ Jahre gestalteten sich schwierig: ich versuchte erst morgens vor der Arbeit zu schreiben, dann abends, dann am Wochenende. Am besten half es, Urlaub zu nehmen, aber der war natürlich begrenzt. Insgesamt habe ich viele Feiertage und Wochenenden der Dissertation geopfert, und kann jedem, der überlegt die Dissertation neben der Arbeit fertig zu stellen nur empfehlen, das vorher zu tun.

Rückblickend bin ich froh, dass ich angefangen habe zu promovieren. Ich bin froh, dass ich mein Thema gewechselt habe. Ich bin dankbar, dass ich mich so lange Zeit intensiv mit einem für mich spannenden Thema beschäftigen konnte. Und ich bin froh, dass ich die Dissertation am Ende noch fertig gestellt habe. 

\newpage
\chapter*{Danksagung}
\pdfbookmark[1]{Danksagung}{danksagung}
Diese Dissertation hätte nicht entstehen können, hätten mich nicht zahlreiche Leute auf dem Weg dahin unterstützt. Ich alleine hätte das Durchhaltevermögen nicht gehabt. Daher bedanke ich mich bei allen, die mich auf dem Weg begleitet haben und entschuldige mich bei denen, die ich hier nicht erwähne.

Mein Dank gilt meiner Familie. Ohne sie hätte ich wahrscheinlich weder studiert noch promoviert. Sie verloren nie den Glauben daran, dass ich die Promotion irgendwann noch abschließen würde. 

Mein Dank gilt den Ilmenauer Pimmelbergern, insbesondere meinen Leidensgenossen Philipp und Seb, aber natürlich auch Lena und Christina. Lena gilt ein besonderer Dank. Sie hat mich in den sechs Jahren, in denen wir zusammen waren und zusammen gelebt haben, bei der Promotion unterstützt, auch wenn das bedeutete, dass ich lange Nächte und Wochenenden vor dem Rechner saß und häufig schlecht gelaunt war.

Mein Dank und meine Liebe gelten Katharina, die sich trotz unseres vollen gemeinsamen Lebens nie beschwert hat, wenn ich Wochenenden, Urlaube und Feiertage der Dissertation geopfert habe.

Mein Dank gilt meinem Doktorvater Thomas, der durch sein wertvolles Feedback dafür gesorgt hat, dass die Dissertation qualitativ nun da ist, wo sie ist.

Mein Dank gilt den weiteren Mitarbeitenden und Professor:innen am Institut für Mathematik, die mich auf meinem Weg begleitet haben. Danke an Achim, der Thomas des Öfteren ermuntert hat, sich bei mir zu melden. 

Danke an meine neuen Arbeitskolleg:innen bei der singularIT, insbesondere Mattis und Natalie K., die mich dazu ermuntert haben, die Dissertation zügig fertigzustellen.

\end{otherlanguage}
\endgroup